
% ДЛЯ КА ВВ и КА СВ! КА СВ отличие только во входах на включение выключателя в схеме В используются xlsx VV_SwitchDevice и SV_SwitchDevice
\needspace{3\baselineskip}
\color{uniblue}{\subsection{Автоматическое регулирование}}
\color{black}

\begin{enumerate}[label=\arabic{section}.\arabic{subsection}.\arabic{enumi}, labelsep=4pt, leftmargin=0pt, itemindent=57pt, itemsep=0pt, parsep=5pt] % если что можно крутить itemsep

%%%%%%%%%%%%%%%%%%%%%%%%%%%%%%%%%%%%%%%%%%%%%%%%%%%%%%%%%%%%%%%%%%%%% ОБЩИЕ СВЕДЕНИЯ %%%%%%%%%%%%%%%%%%%%%%%%%%%%%%%%%%%%%%%%%%%%%%%%%%%%%%%%%%%%%%

\item Автоматическое регулирование

В режиме автоматического регулирования устройство автоматически формирует команды управления приводом с целью поддержания уровня напряжения в заданном диапазоне.

\item Выбор регулируемой секции
    \begin{enumerate}[label=\arabic{section}.\arabic{subsection}.\arabic{enumi}.\arabic*, labelsep=4pt, leftmargin=0em, itemindent=65pt, parsep=0pt]
        \item 
        Устройство РПН, установленное на стороне ВН трехобмоточного трансформатора (рисунок 2.2) или трансформатора с расщепленной обмоткой  НН (рисунок 2.1, б), одновременно воздействует на  уровень напряжения на двух подключенных к трансформатору секциях. В данном случае одна из секций определяется как регулируемая, другая как контролируемая.
        \item
        Устройство в автоматическом режиме осуществляет регулирование по уровню напряжения на регулируемой секции. На контролируемой секции, соответственно, проверяется отсутствие условий блокировки автоматического регулирования.
        \item
        В устройстве предусмотрен следующий механизм определения регулируемой и контролируемой секции. Предусмотрено два режима (определяется состоянием виртуального ключа <<Ввод автоматического выбора регулируемой секции>>):
            \begin{itemize}
                \item автоматический;
                \item оперативный.
            \end{itemize}
        \item
        В автоматическом режиме выбор регулируемой секции осуществляется по состоянию блок-контактов вводных выключателей секций и положению виртуального ключа <<Ввод режима приоритета по 2 секции>>. Если подключена только одна из двух секций, она автоматически становится регулируемой.  Если подключены обе секции выбор регулируемой секции осуществляется оперативно виртуальным ключом <<Ввод режима приоритета по 2 секции>>. Подключенная секция, которая не является регулируемой, определяется термином - контролируемая.В автоматическом режиме выбор регулируемой секции осуществляется по состоянию блок-контактов вводных выключателей секций и положению виртуального ключа <<Ввод режима приоритета по 2 секции>>. Если подключена только одна из двух секций, она автоматически становится регулируемой.  Если подключены обе секции выбор регулируемой секции осуществляется оперативно виртуальным ключом <<Ввод режима приоритета по 2 секции>>. Подключенная секция, которая не является регулируемой, определяется термином - контролируемая.
        \item
        В оперативном режиме выбор регулируемой секции осуществляется с помощью виртуального ключа <<Ввод режима приоритета по 2 секции>>. При выведенном состоянии данного виртуального ключа, регулируемой является 1 секция, при введенном - 2-я. Если регулируемая секция отключена от трансформатора формируется сигнал <<Нет связи с регулируемыми шинами>>. Автоматическое регулирование при этом блокируется.
        \item 
        Дискретные сигналы <<Секция 1 включена>>, <<Секция 1 отключена>>, <<Секция 2 включена>>, <<Секция 2 отключена>>  могут поступать в логику устройства посредством дискретных входов устройства или GOOSE-сообщения.
    \end{enumerate}

\item Зона нечувствительности
    \begin{enumerate}[label=\arabic{section}.\arabic{subsection}.\arabic{enumi}.\arabic*, labelsep=4pt, leftmargin=0em, itemindent=65pt, parsep=0pt]
        \item Требуемый диапазон поддержания напряжения на шинах (<<Зона нечувствительности>>) определяется уставками:
            \begin{itemize}
                \item центр зоны нечувствительности (или напряжение поддержания) Uпод;
                \item ширина зоны нечувствительности dUпод.
            \end{itemize}
        \item 
        Оперативное изменение значения напряжения поддержания в зависимости от суточного графика нагрузки возможно переключением между четырьмя группами уставок. 
        \item
        Значение напряжения поддержания может меняться автоматически в зависимости от тока нагрузки в режиме встречного регулирования. 
        \end{enumerate}
    
\item Встречное регулирование
    \begin{enumerate}[label=\arabic{section}.\arabic{subsection}.\arabic{enumi}.\arabic*, labelsep=4pt, leftmargin=0em, itemindent=65pt, parsep=0pt]
        \item
        Встречное регулирование – регулирование напряжения на шинах удаленного потребителя с учетом расчетного падения напряжения в сети в зависимости от тока нагрузки. Для корректного вычисления величины падения напряжения необходимо соблюдать полярность подводимых к устройству цепей переменного тока.
        \item 
        Падение напряжения в распределительной сети определяется по формуле:

        \begin{equation*}
        \Delta \overline{U} = \left( \overline{I}{\textsl{св}} - \overline{I}{\textsl{вв}} \right) \times \overline{Z}{\textsl{сети}}
        \end{equation*}
        \begin{eqexpl}[25mm]
            \item{$Z_{\textsl{сети}}$} комплексное значение сопротивление сети (линии); определяется по задаваемым параметрам активного и реактивного сопротивления сети; если встречное регулирование не предусматривается, значения данных параметров должны быть равными 0;
            \item{$\overline{I}_{\textsl{вв}}$}комплексное значение тока ввода секции;
            \item{$I_{\textsl{св}}$}комплексное значение тока секционного выключателя.
        \end{eqexpl}
        \item
        Значение напряжения на шинах потребителя определяется по формуле:
        \begin{equation*}
        U{\textsl{потр}} = \mid \overline{U}{\textsl{тек}} - \Delta \overline{U} \mid 
        \end{equation*}

        \begin{eqexpl}[25mm]
            \item{$\overline{U}_{\textsl{тек}}$} комплексное значение напряжения на шинах регулируемой секции.
        \end{eqexpl}

    \end{enumerate}


\item Процесс автоматического регулирования 
    \begin{enumerate}[label=\arabic{section}.\arabic{subsection}.\arabic{enumi}.\arabic*, labelsep=4pt, leftmargin=0em, itemindent=65pt, parsep=0pt]
        \item
        В автоматическом режиме устройство непрерывно отслеживает уровень напряжения на регулируемой секции. В устройстве предусмотрены специальные ИО, которые при выходе напряжения за пределы зоны нечувствительности формируют следующие сигналы:
            \begin{itemize}
                \item напряжение ниже границы зоны нечувствительности;
                \item напряжение выше границы зоны нечувствительности;
                \item напряжение значительно выше границы зоны нечувствительности.
            \end{itemize}
        \item
        В зависимости от возникающего режима и при отсутствии блокирующих сигналов логика формирования команд в автоматическом режиме регулирования с определенной задержкой формирует сигналы управления приводом.

    \end{enumerate}

\item Уставки и параметры функции



\small
\begin{longtable}{|>{\centering\arraybackslash}m{5.3cm}|>{\centering\arraybackslash}m{3.3cm}|>{\centering\arraybackslash}m{4.2cm}|>{\centering\arraybackslash}m{1.8cm}|>{\centering\arraybackslash}m{1cm}|}
\caption{Параметры для настройки функции <<Автоматическое регулирование>>\hfill\vspace{-0.5\baselineskip}}\label{mtzvn:tbl1}\\ 
\hline
\rowcolor{gray!30}
Параметр (Параметр на ИЧМ) & Условное обозначение на схеме & Значение/ Диапазон & Единица измерения & Шаг \\ 
\hline
\endfirsthead
%\caption{\emph{Продолжение\hfill\vspace{-0.5\baselineskip}}} \\ % Отступ обозначения таблицы от таблицы и выравнивание надписи слева
\caption*{\hspace{3pt}\emph{Продолжение таблицы \ref{mtzvn:tbl1}\hfill\vspace{-0.5\baselineskip}}} \\ % сделано по ГОСТ 2.105 п.6.8.7
\hline
\rowcolor{gray!30}
Параметр (Параметр на ИЧМ) & Условное обозначение на схеме & Значение/ Диапазон & Единица измерения & Шаг \\ 
%\hline
\endhead
%\hline
\endfoot
%\hline
\endlastfoot
%==+t1*ATCC1|_TTCC2> АРН

%===t1
%==+t1*AVTR1|_TTCC2> Выбор ТН

%===t1

\end{longtable}
\normalsize

% % ФОРМУЛА №2
% \begin{equation*}
% k_{AM,n}= \frac{I_{\textsl{перв,n}}\cdot  I_{\textsl{ном.терм,n}} }{I_{\textsl{баз,n}}\cdot I_{\textsl{втор,n}}\cdot k_{\textsl{сх,n}}},
% \end{equation*}
% \begin{equation*}
% I_{\textsl{баз,n}}=\frac{S_{\textsl{баз}}}{\sqrt{3} \cdot U_{\textsl{баз,n}} },
% \end{equation*}
% \begin{eqexpl}[25mm]
% \item{$n$} индекс, обозначающий сторону (плечо) ДЗТ;
% \item{$I_{\textsl{перв,n}}$}номинальный первичный ток ТТ стороны $n$;
% \item{$I_{\textsl{втор,n}}$}номинальный вторичный ток ТТ стороны $n$;
% \item{$I_{\textsl{ном.терм,n}}$}номинальное значение тока аналоговых входов устройства, подключенных к ТТ стороны $n$;
% \item{$k_{\textsl{сх,n}}$}коэффициент, учитывающий схему соединения обмоток ТТ стороны~$n$. Данный параметр принимается равным 1, когда ТТ стороны $n$ собраны в звезду, или равным $\sqrt{3}$, когда ТТ стороны $n$ собраны в треугольник. Коэффициенты схемы соединения ТТ для каждой из сторон задается уставками <<KсхемСт1>>, <<KсхемСт2>>, <<KсхемСт3>>;
% \item{$S_{\textsl{баз}}$}базисная мощность, равная номинальной мощности силового трансформатора;
% \item{$U_{\textsl{баз,n}}$}базисное напряжение, равное номинальному напряжению обмотки трансформатора соответствующей стороне $n$.
% \end{eqexpl}

% % ФОРМУЛА №1
% \begin{equation*}
% \begin{bmatrix}
% I_{A,n*}\\ 
% I_{B,n*}\\ 
% I_{C,n*}
% \end{bmatrix}= k_{AM,n}\times \begin{bmatrix}
% M_{11,n} & M_{12,n} & M_{13,n}\\ 
% M_{21,n} & M_{22,n} & M_{23,n}\\ 
% M_{31,n} & M_{32,n} & M_{33,n}
% \end{bmatrix}\times \begin{bmatrix}
% I_{A,n}\\ 
% I_{B,n}\\ 
% I_{C,n}
% \end{bmatrix},
% \end{equation*}
% \begin{eqexpl}[25mm]
% \item{$n$}индекс, обозначающий сторону (плечо) ДЗТ;
% \item{$k_{AM,n}$}коэффициент амплитудного выравнивания;
% \item{$I_{A,n},I_{B,n},I_{C,n}$}измеренные фазные токи стороны $n$;
% \item{$M_{11,n}\dots M_{33,n} $}коэффициенты матрицы компенсации фазового сдвига в соответствие  с группой соединения обмоток трансформатора.
% \end{eqexpl}


%%%%%%%%%%%%%%%%%%%%%%%%%%%%%%%%%%%%%%%%%%%%%%%%%%%%%%%%%%% КОНЕЦ Выкатной элемент (ВЭ) %%%%%%%%%%%%%%%%%%%%%%%%%%%%%%%%%%%%%%%%%%%%%%%%%%%%%%%%%%


\end{enumerate}